\chapter{Conclusion and Future Directions}\label{chap:conclusion}

We have established that knot Floer homology detects the granny knot $T_{2,3}\# T_{2,3}$ as a bigraded vector space over $\mathbb{F} = \mathbb{Z}/2\mathbb{Z}$, providing the first detection result for a composite knot. Along the way, we developed the \texttt{autofolder} algorithm — a complete, algorithmic realization of the folding-automaton framework theorized by Ko, Los, and Song — and a partial detection result for the $(2,1)$-cable of the trefoil. In this concluding chapter, we summarize the principal results, discuss their broader implications, and identify directions for future research.

\section{Summary of Results}

The principal results of this dissertation are the following.

\begin{theorem}[Theorem \ref{thm:main_detection}]
Let $K \subset S^3$ be a knot. If $\widehat{HFK}(K; \mathbb{F}) \cong \widehat{HFK}(T_{2,3}\#T_{2,3}; \mathbb{F})$ as bigraded vector spaces, then $K$ is isotopic to $T_{2,3}\# T_{2,3}$.
\end{theorem}

\begin{theorem}[Theorem \ref{thm:autofolder_main}]
The \texttt{autofolder} algorithm provides a complete, finite classification of all standard train tracks and their valid elementary folds for any specified stratum of pseudo-Anosov braids on the $n$-punctured disk.
\end{theorem}

\begin{theorem}[Theorem \ref{thm:cable_detection_partial}]
Let $K \subset S^3$ be a knot with $\widehat{HFK}(K) \cong \widehat{HFK}(C_{2,1}(T_{2,3}))$ as bigraded vector spaces. Then either $K \cong C_{2,1}(T_{2,3})$, or $K$ is hyperbolic and the disk braid $\beta \in \mathrm{Mod}(D_6)$ associated to the monodromy of $K$ via the Birman-Hilden correspondence (capping off to $\Sigma_2^0$, descent to $S^2_{0,6}$, blowup at the canonical fixed point) belongs to one of finitely many parametric families enumerated in Appendix \ref{app:cable_residual_families}. In particular, $\widehat{HFK}$ detects $C_{2,1}(T_{2,3})$ among non-hyperbolic knots.
\end{theorem}

\begin{corollary}[Corollary \ref{cor:chirality}]
Knot Floer homology distinguishes the homochiral connect-sum of trefoils ($T_{2,3}\#T_{2,3}$) from the heterochiral connect-sum ($T_{2,3}\#\overline{T}_{2,3}$) and from the mirror granny knot ($\overline{T}_{2,3}\#\overline{T}_{2,3}$).
\end{corollary}

The proof strategy throughout combines three threads. First, the genus and fiberedness detection results of Ozsv\'{a}th-Szab\'{o}, Ni, and Juh\'{a}sz translate the topological classification problem into a problem about surface monodromies. Second, the Birman-Hilden correspondence descends the genus-two monodromy to a 6-braid on the disk, where train track methods become applicable. Third, the explicit machinery developed in this dissertation — the \texttt{autofolder}, tight splitting, the Trace Lemma, the Absorption Obstruction, and the homological filtering pipeline — exhaustively eliminates candidate mapping classes within each topological stratum.

\section{Direct Consequences}

\subsection{Almost L-space rigidity}

The granny knot is the unique composite almost L-space knot, by a theorem of Binns \cite{Binns2023}. Combined with our Theorem \ref{thm:main_detection}, this establishes:

\begin{corollary}
Knot Floer homology detects the unique composite almost L-space knot.
\end{corollary}

For the cable, our partial result already establishes detection within the non-hyperbolic class:

\begin{corollary}
Among non-hyperbolic knots in $S^3$, the cable $C_{2,1}(T_{2,3})$ is uniquely identified by its bigraded knot Floer homology.
\end{corollary}

These results contribute to the program of using $\widehat{HFK}$ as a fingerprint for distinguished knot classes, complementing Reinoso's classification of genus-two L-space knots \sidecite{Reinoso2023}, from which $\widehat{HFK}$-detection of $T(2,5)$ follows, and Baldwin-Hu-Sivek's Khovanov-homology detection of $T(2,5)$ \sidecite{baldwin2025khovanov}.

\subsection{The autofolder as a stand-alone contribution}

The \texttt{autofolder} algorithm, while developed here in service of Theorem \ref{thm:main_detection}, is independently useful. Los \sidecite{LosAlgorithm} explicitly called for an algorithmic realization of the folding-automaton machinery to make the state space of pseudo-Anosov mapping classes tractable; the \texttt{autofolder} provides this. Beyond the specific strata analyzed in Chapters \ref{chap:stratoriented} and \ref{chap:stratatraintrackanalysis}, the algorithm operates on arbitrary strata of the marked disk and outputs the complete folding automaton.

We anticipate that the algorithm will be a useful tool for future detection problems on genus-two and higher fibered knots, for the explicit construction of pseudo-Anosov braids of small dilatation, and for the experimental investigation of singularity strata that have not yet been analyzed by hand.

\section{Future Directions}

\subsection{Closing the cable detection gap}

The most immediate open problem left by this dissertation is closing the gap in Theorem \ref{thm:cable_detection_partial}: eliminating the residual finite list of candidate braid families that pass both the homology-sphere filter and the Alexander polynomial filter at non-negative integer parameter values. Three avenues are discussed in Section \ref{sec:cable_residual_discussion}:

\begin{itemize}
    \item Refining the homological filter from the characteristic polynomial of $A_{\text{red}}$ to the full $\mathbb{Z}$-conjugacy class within $\mathrm{Sp}(4, \mathbb{Z})$, capturing the symplectic structure;
    \item Verifying pseudo-Anosov-ness of each candidate braid at the surviving parameter values, since some parameter values may correspond to reducible or periodic mapping classes outside the scope of the topological enumeration;
    \item Comparing the bigraded structure of $\widehat{HFK}$ rather than just the Alexander polynomial, which captures Maslov-grading information beyond what the Euler characteristic detects.
\end{itemize}

Any one of these avenues, if it successfully eliminates all residual candidates, would close the gap and yield a complete cable detection theorem.

\subsection{Detection beyond genus two}

The strategy of this dissertation --- Ni-Juh\'{a}sz reduction to surface monodromy, Birman-Hilden descent to a braid on the disk, and exhaustive train-track elimination of candidate mapping classes --- extends to higher-genus fibered knots only under an important structural constraint.

In genus two, the hyperelliptic involution $\iota$ is central in the mapping class group $\mathrm{Mod}(\Sigma_2)$, so \emph{every} monodromy of a genus-two fibered knot automatically commutes with $\iota$ and therefore descends to a $(2g+2) = 6$-braid on the disk via the Birman-Hilden correspondence. For genus $g \geq 3$, the picture changes: the hyperelliptic mapping class group
\[
\mathcal{H}_g \;=\; C_{\mathrm{Mod}(\Sigma_g)}(\iota)
\]
is a \emph{proper} subgroup of $\mathrm{Mod}(\Sigma_g)$, and Birman-Hilden descent applies only to monodromies in $\mathcal{H}_g$ --- the so-called \textbf{hyperelliptic} (or \textbf{symmetric}) monodromies. Consequently, the framework developed here generalizes naturally to fibered knots of genus $g \geq 3$ \emph{whose monodromy is hyperelliptic}, but not directly to general genus-$g$ fibered knots. For non-hyperelliptic monodromies, an analogous theory would need to operate on closed-surface train tracks directly (see the discussion of \emph{Generalizing the autofolder} below).

Even within the hyperelliptic class, the principal computational obstruction in genus three is the rapid growth in complexity of the train-track state space: the Euler-Poincar\'{e} stratum analysis admits substantially more strata, and the \texttt{autofolder}'s output graphs grow correspondingly. Natural genus-three targets include the torus knot $T_{2,7}$ (whose torus-knot monodromy is hyperelliptic) and the composite knot $T_{2,3} \# T_{2,5}$ (the unique composite genus-three fibered knot of the form $T_{2,p} \# T_{2,q}$); both are amenable to the framework developed here, modulo computational scaling and the verification that the relevant monodromies indeed lie in $\mathcal{H}_3$.

\subsection{Prime detection for satellite operators}

A natural generalization of the cable detection result is the question: for which satellite patterns $P$ is $P(T_{2,3})$ detected by $\widehat{HFK}$? The framework of Schubert's formula combined with the \texttt{autofolder} should apply to a wide range of cabling, doubling, and Mazur-type patterns. We expect that the same machinery, with the homological filter retargeted appropriately, will yield detection results (or partial-detection results in the spirit of Theorem \ref{thm:cable_detection_partial}) for many of the standard pattern families.

\subsection{Khovanov detection of the granny knot}

The detection of the granny knot by knot Floer homology raises the natural question of whether Khovanov homology detects $T_{2,3}\# T_{2,3}$ as well. Khovanov homology is currently known to detect the unknot, the two trefoils, the figure-eight, the cinquefoils $T(\pm 2,5)$, $5_2$, and the Whitehead doubles $\mathrm{Wh}^\pm(T_{2,3},2)$, but is not known to detect any composite knot. A Khovanov-side analog of Theorem \ref{thm:main_detection} would be the first such detection result.

We outline the natural strategy here, adapted from the Khovanov detection of the figure-eight knot by Baldwin, Dowlin, Levine, Lidman, and Sazdanovic \sidecite{BaldwinDowlinLevineLidmanSazdanovic2021}. The granny knot has the following relevant invariants:
\begin{itemize}
    \item $\dim_{\mathbb{Q}} \widetilde{Kh}(T_{2,3}\# T_{2,3}) = 9$, since reduced Khovanov homology is multiplicative under connect sum and $\dim_{\mathbb{Q}} \widetilde{Kh}(T_{2,3}) = 3$;
    \item $T_{2,3}\# T_{2,3}$ is thin (a connect sum of thin knots is thin);
    \item $T_{2,3}\# T_{2,3}$ is strongly quasipositive, and $\tau(T_{2,3}\# T_{2,3}) = g(T_{2,3}\# T_{2,3}) = 2$.
\end{itemize}

Let $K \subset S^3$ be a knot whose reduced Khovanov homology over $\mathbb{Q}$ is isomorphic to $\widetilde{Kh}(T_{2,3}\# T_{2,3};\mathbb{Q})$ as a bigraded vector space. The proposed argument proceeds in three steps:

\begin{enumerate}
    \item \emph{Khovanov thinness.} The bigraded isomorphism implies $K$ is thin (Khovanov-thinness is a property of the bigraded structure).

    \item \emph{Spectral sequence collapse.} Dowlin's spectral sequence \sidecite{Dowlin2024} from reduced Khovanov homology to $\widehat{HFK}$ degenerates at the $E_2$-page for thin knots, yielding
    \[
    \widehat{HFK}(K;\mathbb{Q}) \cong \widehat{HFK}(T_{2,3}\#T_{2,3};\mathbb{Q})
    \]
    as bigraded vector spaces.

    \item \emph{$\widehat{HFK}$-detection.} Theorem \ref{thm:main_detection} (extended from $\mathbb{F}_2$ to $\mathbb{Q}$ coefficients) concludes $K \cong T_{2,3}\# T_{2,3}$.
\end{enumerate}

The principal subtlety in this strategy is the extension of Theorem \ref{thm:main_detection} from $\mathbb{F}_2$ to $\mathbb{Q}$ coefficients, which requires verifying that the fixed-point obstructions used in the proof remain valid in characteristic zero. We do not pursue this in the present work.

\subsection{$\widehat{HFK}$ and Khovanov homology on the Wang family}

The square knot $T_{2,3} \# \overline{T}_{2,3}$ is a band sum of the split link $T_{2,3} \sqcup \overline{T}_{2,3}$ along a trivial band. By a construction of Hedden-Watson \sidecite{HeddenWatson2018}, generalized by Wang \sidecite{Wang2020}, adding $n$ half-twists to such a band produces an infinite family $\{K_n\}_{n\in\mathbb{Z}}$ of distinct knots, all sharing the same knot Floer homology (and the same instanton knot Floer homology), of which $K_0$ is the square knot. The Khovanov homology of the knots $K_n$ are mutually distinct.

This picture suggests two complementary detection questions:

\begin{enumerate}
    \item \emph{Does Khovanov homology distinguish each $K_n$ from all knots outside the Wang family?} The Wang family is precisely the family of obstructions to $\widehat{HFK}$-detection of the square knot, so a positive answer for the square knot ($n=0$) would yield the first Khovanov detection result for a composite knot. The strategy outlined for the granny knot in the previous subsection does not directly apply to the square knot, since $\tau(T_{2,3} \# \overline{T}_{2,3}) = 0 \neq g(T_{2,3} \# \overline{T}_{2,3}) = 2$, and the strong-quasipositive simplifications used in the granny case fail.

    \item \emph{Does Khovanov homology detect membership in the Wang family?} That is: if $\widetilde{Kh}(K) \cong \widetilde{Kh}(K_n)$ for some $n$, must $K = K_m$ for some $m$? An affirmative answer would establish a striking complementarity between $\widehat{HFK}$ and Khovanov homology: $\widehat{HFK}$ identifies the family as an isomorphism class, while Khovanov resolves membership within the family. To our knowledge, no knot invariant is currently known to detect a non-trivial infinite family of knots in this sense, so such a result would be a first.
\end{enumerate}

The first question is the natural Khovanov analog of the square knot detection problem, which lies beyond the scope of $\widehat{HFK}$. The second is a structural question about the relationship between $\widehat{HFK}$ and Khovanov homology that is naturally raised by, but does not obviously follow from, the techniques developed in this dissertation.

\subsection{Generalizing the autofolder}

The \texttt{autofolder} as currently implemented enumerates standard train tracks on the marked disk. Two natural extensions are:

\begin{itemize}
    \item \textbf{Closed surfaces of higher genus.}  Train track theory on closed surfaces (rather than on the disk) is well-developed, but the algorithmic enumeration is more delicate due to the absence of a canonical embedding.  An autofolder operating directly on $\Sigma_g$ would bypass the Birman-Hilden descent step and could enable detection results for fibered knots in 3-manifolds other than $S^3$.
    \item \textbf{Activating the edge data.} The current \texttt{autofolder} outputs the vertex set of the folding automaton; the edge set, encoding standardizing braids and transition matrices, is computed by the package but not utilized in our obstruction. The dilatation-enumeration applications of this edge data are discussed in the introduction (Section \ref{subsec:dilatation_enumeration_intro}); other potential uses include the study of finer dynamical invariants of pseudo-Anosov mapping classes within each stratum.
\end{itemize}

\section{Closing Remarks}

The detection of the granny knot by knot Floer homology demonstrates that the bigraded structure of $\widehat{HFK}$ is sufficiently rich to identify a knot uniquely even in the presence of essential tori in its complement, where geometric obstructions of the hyperbolic kind are unavailable. The combination of monodromy descent via Birman-Hilden, train-track classification via the \texttt{autofolder}, and homological filtering via the characteristic polynomial of the lifted action provides a template for future detection results in low-dimensional topology.

We hope that the methods developed in this dissertation — particularly the algorithmic infrastructure for navigating the geometry of mapping class groups — will serve as a springboard for further work on the rigidity of knot invariants and the classification of pseudo-Anosov mapping classes.